\paragraph{A counterexample to the full flat, locally finitely presented assertion.}

\begin{theorem}\label{thm:DE-E2-counterexample}
\status{Proved; direct all-tests argument independently audited by the parent and reduction-descent agents.}
For every prime $p$ there is an $\Fp$-algebra $R$ and a finite free morphism $X\to\Spec R$ of rank $p$ for which the restriction functor
\[
 \{\text{crystals of $\OO$-modules on $(X/R)_{\rm inf}$}\}
 \longrightarrow
 \{\text{ordinary stratified $\OO_X$-modules relative to $R$}\}
\]
is not faithful.  The failure already occurs among automorphisms of the rank-one locally free structure crystal.  Thus Part~E.2, with precisely its stated hypotheses, has a negative answer.
\end{theorem}
\begin{proof}
Put
\[
 R_0=\Fp[t],\qquad I=\Fp[t,t^{-1}]/\Fp[t],\qquad
 R=R_0\ltimes I.
\]
Here the last expression is the square-zero extension with multiplication
$(a,u)(b,v)=(ab,av+bu)$.  We regard $t$ as $(t,0)$ and $I$ as the ideal $0\oplus I$ of $R$.  For $n\geq0$ set
\[
 s_n=1+p+\cdots+p^{n-1}=\frac{p^n-1}{p-1},\qquad
 b_n=(0,[t^{-s_n}])\in I,
\]
where $s_0=0$, so that $b_0=0$.  These elements satisfy
\begin{equation}\label{eq:DE-E2-Pruefer-recurrence}
 t^{p^n}b_{n+1}=b_n\quad(n\geq0),\qquad
 tb_1=0,\qquad b_1\ne0.
\end{equation}
Let
\[
 f=x^p-tx,\qquad A=R[x]/(f),\qquad X=\Spec A.
\]
The polynomial $f$ is monic.  Hence $A$ is a free $R$-module with basis $1,x,\ldots,x^{p-1}$; in particular $X/R$ is flat and finitely presented.  Moreover $f$ is a non-zero-divisor in $R[x]$, and every fibre is a zero-dimensional hypersurface over a field, hence a complete intersection and Cohen--Macaulay.  The example does not rely on a failure of the Cohen--Macaulay property of the fibres.

We construct a nonzero endomorphism $\delta$ of the structure crystal whose evaluation on $X$ is zero.  Let $U\hookrightarrow T$ be any infinitesimal test object, with nilpotent ideal $J$.  Locally on $T$, choose an integer $n$ such that $J^{[p^n]}=0$ and a lift $z\in\OO_T$ of the function $x|_U$.  Define
\begin{equation}\label{eq:DE-E2-phantom-value}
                     \delta_T=b_nz^{p^n}\in\OO_T.
\end{equation}
Changing $z$ by an element of $J$ does not change this function, because $T$ has characteristic $p$.  Also $z^p-tz\in J$, and therefore
\[
 b_{n+1}z^{p^{n+1}}-b_nz^{p^n}
       =b_{n+1}(z^p-tz)^{p^n}=0
\]
by \eqref{eq:DE-E2-Pruefer-recurrence}.  The definition is consequently independent of every sufficiently large $n$.  These two observations show that the local functions glue, including when only local nilpotence bounds have been chosen.

For a morphism of infinitesimal test objects, choose a common sufficiently large $n$ and pull back a local lift $z$.  Formula~\eqref{eq:DE-E2-phantom-value} is then preserved by pullback.  Thus the functions $\delta_T$ define a natural $\OO$-linear endomorphism of the structure crystal.  This construction works on all test objects, without a finiteness or quasi-coherence restriction on their nilpotent ideals.  It also works for the big infinitesimal site if that convention is adopted.

On the test object $X\hookrightarrow X$ one may take $n=0$, giving
$\delta_X=b_0x=0$.  On the other hand, take the square-zero test
\[
             X\hookrightarrow T=\Spec B,\qquad B=R[x]/(f^2).
\]
Here $J=(f)/(f^2)$ satisfies $J^2=0$, so $n=1$ is allowed for every $p$.  Formula~\eqref{eq:DE-E2-phantom-value} gives
\[
                  \delta_T=b_1x^p=b_1f\ne0.
\]
The equality uses $tb_1=0$.  Nonvanishing follows from the monicity of $f^2$: the classes of $1,x,\ldots,x^{2p-1}$ form an $R$-basis of $B$, and the coefficient of $x^p$ is $b_1\ne0$.  In fact this test $T$ is itself finite free over $R$.

Finally, $I^2=0$ implies $\delta_T^2=0$ for every $T$.  Hence $1+\delta$ is an automorphism of the rank-one structure crystal, with inverse $1-\delta$, and its restriction to the ordinary stratified module $\OO_X$ is the identity.  The restriction functor is not faithful, proving the assertion.
\end{proof}


\begin{corollary}[Failure over a regular Noetherian base with arbitrary quasi-coherent coefficients]\label{cor:DE-E2-Noetherian-counterexample}
\status{Proved; independently audited by the parent and reduction-descent agents.}
The restriction functor of Theorem~\ref{thm:DE-E2-counterexample} can fail to be faithful even when the base is the regular Noetherian ring $R=\Fp[t]$ and $X/R$ is finite free with complete-intersection fibres.  In this version both crystals are quasi-coherent; one is the rank-one structure crystal and the other has arbitrary, noncoherent coefficients.
\end{corollary}
\begin{proof}
Use $R=\Fp[t]$, the $R$-module $I=R[t^{-1}]/R$, and $A=R[x]/(x^p-tx)$.  We use no square-zero extension of the base.  Besides the structure crystal, consider the constant-base coefficient crystal
\[
                 \mathcal E_T=\OO_T\otimes_R I.
\]
Its crystal transition isomorphisms are ordinary tensor associativity, so no flatness assumption on $I$ is required.  Each value is a quasi-coherent module.

With $s_n$ and $b_n=[t^{-s_n}]\in I$ as before, the sections
\[
                    z^{p^n}\otimes b_n\in\mathcal E_T
\]
define a morphism $\OO_{\rm inf}\to\mathcal E$.  The lift-independence, compatibility in $n$, gluing, and naturality on all test objects are exactly the two identities proved in Theorem~\ref{thm:DE-E2-counterexample}, now interpreted in a module rather than a square-zero ideal.  Its value on $X$ is zero, whereas its value on $B=R[x]/(x^p-tx)^2$ is
\[
                         x^p\otimes b_1\ne0.
\]
Indeed, $B\otimes_R I$ is the direct sum of $2p$ copies of $I$ in the monic basis, and its $x^p$ coefficient is $b_1\ne0$.  This proves nonfaithfulness.  Equivalently, the upper-triangular automorphism of $\OO_{\rm inf}\oplus\mathcal E$ with this off-diagonal entry is a nonidentity automorphism invisible to restriction.

The module $I$ is not finitely generated over $R$.  Since $A$ is finite faithfully flat and free over $R$, the module $A\otimes_R I$ is not finitely generated over $A$: otherwise it would be finitely generated over $R$, as would its direct summand $I$.  Thus this example does not fall under the manuscript's coherent-coefficient proviso.
\end{proof}

\begin{remark}[The exact role of the base hypothesis]\label{rem:DE-E2-Noetherian-scope}
The ring $R$ in Theorem~\ref{thm:DE-E2-counterexample} is non-Noetherian: its square-zero ideal $I$ is not finitely generated, since finitely many Laurent principal parts have a bounded pole order whereas $I$ contains every $t^{-m}$.  The supplied transcription at cote~6, scan~7 adds ``du moins si R loc.noeth., et se bornant aux Modules coh\'erents''.  The counterexample refutes the broader assertion actually requested in Part~E.2, even with finite locally free coefficients, and does not refute that inserted locally Noetherian version.  Corollary~\ref{cor:DE-E2-Noetherian-counterexample} further refutes the all-module assertion over a regular Noetherian base.  Only the simultaneous locally Noetherian and coherent-coefficient refinement remains undecided here; the coherent characteristic-zero field theorem is a positive special case.  For locally Noetherian bases in characteristic $p$, the pure-relative-Frobenius result gives another positive special case.

There is also a direct description using Proposition~\ref{prop:DE-E2-relative-Frobenius}.  The $n$th Frobenius twist has algebra
\[
 A_n=R[x_n]/(x_n^p-t^{p^n}x_n),
\]
and multiplication by $a_n=b_nx_n$ defines an endomorphism of the unit $F$-divided system, since
\[
 F_{X^{(n)}/R}^{\#}(a_{n+1})
       =b_{n+1}x_n^p=b_nx_n=a_n.
\]
Here $a_0=0$ but $a_1\ne0$.  This is an infinitely compatible failure, not merely an ineffective datum for one Frobenius map.  The preceding direct proof is independent of the $F$-divided comparison theorem.
\end{remark}
